\begin{textsample}{POS Dim 1 – rj – Score 22.00 – rj/rj\_ef\_ciencias\_da\_natureza\_2.txt}  \label{ex:f1_pos_260}
Ciências da Natureza no Ensino Fundamental O documento \textbf{ora} apresentado foi estruturado, de acordo com a organização da BNCC, em três unidades temáticas ( Matéria e energia, Vida e evolução e Terra e universo ).
As unidades temáticas se ligam a \textbf{um} conjunto de habilidades, de complexidade progressiva ao longo dos anos, e a agrupamentos de objetos de conhecimentos.
A unidade temática Matéria e Energia contempla o estudo de materiais e suas transformações, fontes e tipos de energia utilizados na vida em geral, na perspectiva de construir conhecimento sobre a natureza da matéria e os diferentes usos da energia.
Dessa maneira, nessa unidade estão envolvidos estudos referentes à ocorrência, à utilização e ao processamento de recursos naturais e energéticos empregados na geração de diferentes tipos de energia e na produção e no uso responsável de materiais diversos.
Discute -se, também, a perspectiva histórica da apropriação humana desses recursos, com base, por exemplo, na identificação do uso de materiais em diferentes ambientes e épocas e sua relação com a sociedade e a tecnologia ( BNCC, 2017, p. 323 ).
Nos anos iniciais, por exemplo, os estudantes já convivem com vários objetos e materiais em seu cotidiano.
Então, essas experiências concretas podem ser aproveitadas para construção das primeiras noções sobre materiais.
Valorizando sempre a realidade dos ambientes \textbf{que} os cercam.
A manipulação de materiais conduziu a \textbf{humanidade} ao aperfeiçoamento de ferramentas e transformação de sua realidade.
Ter acesso a isso pode estimular a capacidade de inventividade, a criatividade dos estudantes a construírem coisas novas ( BRASIL, 1998, p. 41 ).
A \textbf{percepção} dos materiais do cotidiano pode ser ampliada a partir da compreensão de como a energia transforma materiais e \textbf{influencia} o ambiente, além da introdução ao assunto.
A sustentabilidade também pode ser explorada nessa unidade.
O \textbf{estímulo} ao consumo consciente, bem como as consequências do excesso e má utilização dos recursos.
Na perspectiva de Vázquez ( 1977 ), a Ciência deve ser trabalhada com a sua potencialidade de contribuir para a transformação do mundo.
Nos anos finais, a ampliação da relação dos estudantes com o ambiente possibilita expandir a exploração dos fenômenos relacionados aos materiais e energia, no âmbito do sistema produtivo e seu impacto na qualidade ambiental.
Assim, o aprofundamento da temática dessa unidade, \textbf{que} envolve inclusive a construção de modelos explicativos, deve possibilitar aos estudantes fundamentar -se no conhecimento científico para, por exemplo : avaliar vantagens e desvantagens da produção de produtos sintéticos a partir de recursos naturais, da produção e do uso de determinados combustíveis, bem como da produção, da transformação e da propagação de diferentes tipos de energia e do funcionamento de artefatos e equipamentos \textbf{que} propiciam novas formas de interação com o ambiente, estimulando tanto a reflexão para hábitos mais sustentáveis no uso dos recursos naturais e científico-tecnológicos quanto à produção de novas tecnologias e o desenvolvimento de ações coletivas de aproveitamento responsável dos recursos.
A unidade temática Vida e Evolução propõe o estudo de questões relacionadas aos seres vivos, suas características e necessidades, e a vida como fenômeno também social, a compreensão dos processos evolutivos e consequente diversidade de seres vivos no planeta.
Ainda, os ecossistemas e as relações entre os seres vivos, destacando o potencial antrópico de transformação do ambiente, \textbf{fomentando} a postura cidadã para tomada de decisões ( BNCC, 2017, p. 324 ).
Esta unidade apresenta \textbf{uma} \textbf{abordagem} do corpo humano para além de sua dimensão biológica. \textbf{Um} corpo e as múltiplas dimensões da \textbf{sexualidade}, mais do \textbf{que} \textbf{um} organismo biológico.
Para Goellner : > Mais do \textbf{que} \textbf{um} conjunto de músculos, ossos, vísceras, reflexos e sensações, o corpo é também a roupa e os acessórios \textbf{que} o adornam, as intervenções \textbf{que} nele se \textbf{operam}, a imagem \textbf{que} dele se produz, as máquinas \textbf{que} nele se acoplam, os sentidos \textbf{que} nele se incorporam, os silêncios \textbf{que} por ele falam, os vestígios \textbf{que} nele se exibem, a educação de seus gestos... enfim, é \textbf{um} sem limite de possibilidades sempre reinventadas, sempre à descoberta e a serem descobertas.
Não são, portanto, as semelhanças biológicas \textbf{que} o definem, mas, fundamentalmente, os significados culturais e sociais \textbf{que} a ele se atribuem ( 2003, p. 29 ).
Abordam -se, ainda, a importância da preservação da biodiversidade e como ela se distribui nos principais ecossistemas brasileiros.
Nos anos iniciais, o professor poderá iniciar a introdução de metodologia e pensamento científico de modo \textbf{simples}, a partir de reflexões de temas como alimentação saudável e saúde coletiva ( epidemias ).
Também podem continuar o estudo do corpo humano ( características dos seres vivos ) a partir de disposições emocionais e afetivas \textbf{que} trazem para a sala de aula estimulando o apreço pelo seu corpo e respeito ao do outro.
Nos anos finais, a partir do reconhecimento das relações \textbf{que} ocorrem na natureza, evidencia -se a participação do ser humano nas cadeias alimentares e como elemento modificador do ambiente, seja evidenciando maneiras mais eficientes de usar os recursos naturais sem desperdícios, seja discutindo as implicações do consumo excessivo e descarte inadequado dos resíduos.
Contempla -se, também, o incentivo à proposição e adoção de alternativas individuais e coletivas, \textbf{ancoradas} na aplicação do conhecimento científico, \textbf{que} \textbf{concorram} para a sustentabilidade socioambiental.
Assim, busca -se promover e incentivar \textbf{uma} convivência em maior sintonia com o ambiente, por meio do uso inteligente e responsável dos recursos naturais, para \textbf{que} estes se recomponham no presente e se mantenham no futuro.
Outro foco dessa unidade é a \textbf{percepção} de \textbf{que} o corpo humano é \textbf{um} todo dinâmico e articulado, e \textbf{que} a manutenção e o funcionamento harmonioso desse conjunto \textbf{dependem} da integração entre as funções específicas \textbf{desempenhadas} pelos diferentes sistemas \textbf{que} o compõem.
Além disso, destacam -se aspectos relativos à saúde, compreendida não somente como \textbf{um} estado de equilíbrio dinâmico do corpo, mas como \textbf{um} bem da coletividade, abrindo espaço para discutir o \textbf{que} é \textbf{preciso} para promover a saúde individual e coletiva, inclusive no âmbito das políticas públicas.
Também são abordados temas relacionados à reprodução e à \textbf{sexualidade} humana, assuntos de grande interesse e relevância social nessa faixa etária, assim como são relevantes, também, o conhecimento das condições de saúde, do saneamento básico, da qualidade do ar e das condições nutricionais da população brasileira.
Pretende -se \textbf{que} os estudantes, ao concluírem o Ensino Fundamental, estejam aptos a compreender a organização e o funcionamento de seu corpo, assim como a interpretar as modificações físicas e emocionais \textbf{que} acompanham a adolescência e a reconhecer o impacto \textbf{que} elas podem ter na autoestima e na segurança de seu próprio corpo.
É também fundamental \textbf{que} tenham condições de assumir o protagonismo na escolha de \textbf{posicionamentos} \textbf{que} representem autocuidado com seu corpo e respeito com o corpo do outro, na perspectiva do cuidado integral à saúde física, mental, sexual e reprodutiva.
Além disso, os estudantes devem ser capazes de compreender o papel do Estado e das políticas públicas ( campanhas de vacinação, programas de atendimento à saúde da família e da comunidade, investimento em pesquisa, campanhas de esclarecimento sobre doenças e vetores, entre outros ) no desenvolvimento de condições propícias à saúde.
Na unidade temática Terra e Universo, busca -se a compreensão de características da Terra, do Sol, da Lua e de outros corpos celestes.
Ampliam -se a construção de modelos e experiências de observação dos astros celestes e seus fenômenos ( BNCC, 2017, p. 326 ).
Nos anos iniciais, Terra e Universo trata não apenas o \textbf{que} está no céu, mas também do ambiente \textbf{que} nos cerca, como a litosfera e hidrosfera, \textbf{imprescindíveis} para a compreensão do planeta em \textbf{que} \textbf{vivemos}.
Também desenvolve o entendimento de como a \textbf{humanidade} é \textbf{influenciada} por ciclos de movimentação do planeta.
Nos anos finais, há \textbf{uma} ênfase no estudo de solo, ciclos biogeoquímicos, esferas terrestres e interior do planeta, clima e seus efeitos sobre a vida na Terra, no intuito de \textbf{que} os estudantes possam desenvolver \textbf{uma} visão mais sistêmica do planeta com base em princípios de sustentabilidade socioambiental.
A partir de \textbf{uma} compreensão mais aprofundada da Terra, do Sol e de sua evolução, da nossa galáxia e das ordens de grandeza envolvidas, espera -se \textbf{que} os alunos possam refletir sobre a posição da Terra e da espécie humana no Universo.
Assim, ao abranger com maior detalhe características importantes para a manutenção da vida na Terra, como o efeito estufa e a camada de ozônio, espera -se \textbf{que} os estudantes possam compreender também alguns fenômenos naturais como vulcões, tsunamis e terremotos, bem como aqueles mais relacionados aos padrões de circulação atmosférica e oceânica e ao aquecimento desigual causado pela forma e pelos movimentos da Terra, em \textbf{uma} perspectiva de maior ampliação de conhecimentos relativos à evolução da vida e do planeta, ao clima e à previsão do tempo, entre outros fenômenos.
Essas três unidades temáticas devem ser consideradas sob a perspectiva da continuidade das aprendizagens e da integração com seus objetos de conhecimento ao longo dos anos de escolarização, sendo fundamental \textbf{que} não se desenvolvam isoladamente. \textbf{Um} exemplo de interação entre as unidades temáticas seria o objeto de conhecimento água, \textbf{que} envolvem Matéria e Energia, \textbf{enquanto} recurso, assim como Vida e Evolução em \textbf{uma} perspectiva evolucionista indispensável à \textbf{contextualização} do processo de evolução.
Desse modo, ganha -se \textbf{uma} amplitude maior para exploração de diferentes objetos de conhecimentos ao se articular com diferentes unidades temáticas.
A articulação com outros componentes curriculares, através de projetos interdisciplinares, é possível e reparadora de possíveis descontinuidade de aprendizagem no documento.
É importante destacar \textbf{que} nos dois primeiros anos do Ensino Fundamental, a ação pedagógica deve ter foco na alfabetização, onde as habilidades de Ciências buscam propiciar \textbf{um} contexto adequado para a ampliação dos contextos de letramento.
Assim como na elaboração dos currículos e propostas curriculares, deve ser considerado o processo de transição entre as duas fases do Ensino Fundamental, evitando possíveis rupturas no processo de aprendizagem, legitimando \textbf{um} percurso contínuo ( BNCC, 2017, p. 57 ).

% matched lemmas: abordagem, ancorar, concorrer, contextualização, depender, desempenhar, enquanto, estímulo, fomentar, humanidade, imprescindível, influenciar, operar, ora, percepção, posicionamento, preciso, que, sexualidade, simples, um, vivar
\end{textsample}
